%========================================================
% CHAPTER 6 — Frequency-Dependent Fiber Propagation
%              and Polarization-Mode Dispersion
%========================================================

\chapter{Frequency-Dependent Fiber Propagation and Polarization-Mode Dispersion}
\label{ch:frequency_dependent_fiber_propagation}

%========================================================
% Section 6.1 — Motivation and Modeling Hierarchy
%========================================================

\section{Motivation and Modeling Hierarchy}
\label{sec:pmd_motivation_and_hierarchy}

The effective fiber models developed in
Chapter~\ref{ch:effective_fiber_models} describe polarization propagation
through coherent phase evolution, statistical phase averaging, and local
isotropic depolarization. They provide analytically controlled models for
studying changes in polarization correlations, purity, fidelity,
entanglement, and Bell-inequality violation.

These models assume that each photon is effectively monochromatic and that
the fiber transformation is independent of optical frequency. A lossless
fiber is therefore represented by a single local unitary operator acting on
each polarization qubit, and the two-photon evolution is given by
Eq.~\eqref{eq:local_unitary_fiber_model_ch5}. At this level, deterministic
birefringence produces only a coherent polarization transformation: global
purity and bipartite entanglement are preserved under local unitary evolution
\cite{nielsen_chuang_qcqi}.

\medskip

Photons generated through spontaneous parametric down-conversion are,
however, finite-bandwidth wave packets. Their joint spectral structure is
determined by the pump spectrum and by the phase-matching response of the
nonlinear medium
\cite{grice_walmsley_spdc_spectral_1997}. Different frequency components may
therefore experience different polarization transformations during
propagation.

The local fiber operator must consequently be extended to
%
\begin{equation}
    U_K
    \longrightarrow
    U_K(\omega_K),
    \qquad
    K\in\{A,B\}.
    \label{eq:frequency_dependent_local_unitary}
\end{equation}
%
For a resolved frequency pair, the corresponding two-photon polarization
operator is
%
\begin{equation}
    U_{AB}(\omega_A,\omega_B)
    =
    U_A(\omega_A)
    \otimes
    U_B(\omega_B).
    \label{eq:frequency_dependent_two_arm_unitary}
\end{equation}

Propagation of the complete polarization--frequency state remains unitary in
an ideal lossless fiber. Nevertheless, the frequency dependence of
\(U_{AB}(\omega_A,\omega_B)\) can correlate polarization with frequency and
arrival time. If these degrees of freedom are not resolved by the
measurement, the observable polarization state is obtained through the
partial trace
%
\begin{equation}
    \rho_{AB}^{(\mathrm{pol,out})}
    =
    \operatorname{Tr}_{\omega_A,\omega_B}
    \left[
        \rho_{AB}^{(\mathrm{out})}
    \right].
    \label{eq:frequency_partial_trace}
\end{equation}
%
The reduced state may therefore be mixed even when the complete output state
is pure
\cite{nielsen_chuang_qcqi,poon_law_stochastic_pmd_2008}.

\medskip

The physical origin of this coupling is birefringence. In a birefringent
fiber, two orthogonal polarization eigenmodes experience different
propagation constants and consequently accumulate a relative phase. When
this relative phase varies with frequency, the two modes also acquire
different group delays. Their relative temporal displacement is quantified
by the differential group delay, while its accumulation in a spatially
nonuniform fiber gives rise to polarization-mode dispersion
\cite{foschini_poole_statistical_pmd_1991}.

A photon prepared in a superposition of the two polarization modes can thus
emerge in temporally displaced components. If their arrival time is not
resolved, the resulting loss of temporal overlap reduces the off-diagonal
coherence of the polarization state. For polarization-entangled photon
pairs, this mechanism can reduce Bell-state fidelity, concurrence, and
nonlocal correlations
\cite{brodsky_loss_polarization_entanglement_2011,
antonelli_shtaif_brodsky_pmd_2011}.

Unlike the isotropic depolarizing channel introduced previously, this
degradation is not postulated through a phenomenological noise parameter. It
emerges from frequency-dependent local unitary propagation followed by the
elimination of the unresolved spectral degrees of freedom. Its strength
therefore depends on the photon spectrum and on the birefringent response of
the two fiber arms.

\medskip

The development of this chapter proceeds from the frequency-dependent model
of a single birefringent segment to an ordered concatenation representing a
nonuniform fiber. The effective principal polarization states and
differential group delay are then connected to the reduced two-photon
polarization state obtained after spectral tracing.

The objective is to extend the effective models of the previous chapter with
the minimum physical structure required to describe PMD-induced degradation
of polarization entanglement, while retaining a formulation suitable for
numerical simulation.

%========================================================
% Section 6.2 — Spectral Description of SPDC Photon Pairs
%========================================================

\section{Spectral Description of Finite-Bandwidth SPDC Photon Pairs}
\label{sec:spectral_description_spdc}

The polarization Hilbert space and the encoding
\(|H\rangle\equiv|0\rangle\), \(|V\rangle\equiv|1\rangle\) were introduced
in the preceding chapters. To describe frequency-dependent propagation, the
state space of photon \(K\) is extended to
%
\begin{equation}
    \mathcal{H}_K
    =
    \mathcal{H}_{\omega_K}
    \otimes
    \mathcal{H}_{\mathrm{pol},K},
    \qquad
    K\in\{A,B\},
    \label{eq:single_photon_frequency_polarization_space}
\end{equation}
%
where \(\mathcal{H}_{\omega_K}\) describes the continuous frequency degree
of freedom and \(\mathcal{H}_{\mathrm{pol},K}\) is the polarization-qubit
space. Frequency eigenstates are denoted by \(|\omega_K\rangle\) and satisfy
the continuous normalization
\(\langle\omega_K|\omega_K'\rangle
=\delta(\omega_K-\omega_K')\)
\cite{nielsen_chuang_qcqi,poon_law_stochastic_pmd_2008}.

%--------------------------------------------------------
% Subsection 6.2.1 — Joint Spectral State
%--------------------------------------------------------

\subsection{Joint Spectral State of the Photon Pair}
\label{subsec:two_photon_spdc_spectral_state}

The finite-bandwidth photon pair generated through SPDC is represented by
the joint spectral state
%
\begin{equation}
    |F\rangle_{\omega_A\omega_B}
    =
    \iint
    d\omega_A\,d\omega_B\,
    F(\omega_A,\omega_B)
    |\omega_A,\omega_B\rangle,
    \label{eq:spdc_two_photon_spectral_state}
\end{equation}
%
where \(F(\omega_A,\omega_B)\) is the joint spectral amplitude and
\(|\omega_A,\omega_B\rangle
=|\omega_A\rangle\otimes|\omega_B\rangle\).

The amplitude is normalized according to
%
\begin{equation}
    \iint
    d\omega_A\,d\omega_B\,
    |F(\omega_A,\omega_B)|^2
    =
    1.
    \label{eq:joint_spectral_amplitude_normalization}
\end{equation}
%
The corresponding joint spectral intensity is
%
\begin{equation}
    J(\omega_A,\omega_B)
    =
    |F(\omega_A,\omega_B)|^2,
    \label{eq:joint_spectral_intensity}
\end{equation}
%
and gives the joint probability density of the two photon frequencies.

For an SPDC source, the joint spectral amplitude can be expressed
schematically as
%
\begin{equation}
    F(\omega_A,\omega_B)
    =
    \mathcal{N}\,
    \alpha_p(\omega_A+\omega_B)
    \phi_{\mathrm{pm}}(\omega_A,\omega_B),
    \label{eq:spdc_jsa_pump_phase_matching}
\end{equation}
%
where \(\alpha_p\) is the pump spectral envelope,
\(\phi_{\mathrm{pm}}\) is the phase-matching function of the nonlinear
medium, and \(\mathcal{N}\) is a normalization constant
\cite{grice_walmsley_spdc_spectral_1997}.

The pump envelope constrains the sum frequency, whereas phase matching
selects the efficiently generated frequency pairs. Consequently,
\(J(\omega_A,\omega_B)\) need not factorize into independent single-photon
spectra: the two photons may retain spectral correlations or
anticorrelations.

The finite spectral bandwidth also defines a finite temporal coherence
scale. A differential delay introduced by the fibers can therefore make
different polarization components temporally distinguishable even when
their arrival time is not explicitly measured.

%--------------------------------------------------------
% Subsection 6.2.2 — Combined Input State
%--------------------------------------------------------

\subsection{Combined Polarization--Frequency Input State}
\label{subsec:combined_spdc_polarization_frequency_state}

The present model assumes that, at the source output, the joint spectral
state is independent of the polarization state. For the reference Bell state
used throughout this work, the complete input state is
%
\begin{equation}
    |\Psi_{\mathrm{in}}\rangle
    =
    |F\rangle_{\omega_A\omega_B}
    \otimes
    |\Psi^+\rangle_{\mathrm{pol}}.
    \label{eq:factorized_spdc_frequency_polarization_state}
\end{equation}

Equivalently,
%
\begin{equation}
    |\Psi_{\mathrm{in}}\rangle
    =
    \iint
    d\omega_A\,d\omega_B\,
    F(\omega_A,\omega_B)
    |\omega_A,\omega_B\rangle
    \otimes
    |\Psi^+\rangle_{\mathrm{pol}}.
    \label{eq:complete_spdc_input_state_integral}
\end{equation}

This factorization means that the two polarization alternatives forming the
Bell state share the same joint spectral amplitude. The source therefore
introduces no initial spectral information capable of distinguishing them.
Possible source-induced polarization--frequency correlations are excluded
so that the subsequent loss of polarization coherence can be attributed to
fiber propagation.

%--------------------------------------------------------
% Subsection 6.2.3 — Frequency-Dependent Propagation
%--------------------------------------------------------

\subsection{Frequency-Dependent Propagation and Spectral Trace}
\label{subsec:frequency_dependent_local_propagation}

For a lossless fiber, each frequency component undergoes a unitary
polarization transformation. The complete operator acting on photon \(K\)
can be written as
%
\begin{equation}
    \mathcal{U}_K
    =
    \int d\omega_K\,
    |\omega_K\rangle
    \langle\omega_K|
    \otimes
    U_K(\omega_K),
    \qquad
    K\in\{A,B\}.
    \label{eq:frequency_controlled_fiber_operator}
\end{equation}

This has the structure of a frequency-controlled polarization gate: the
frequency component is preserved, while its polarization undergoes the
corresponding transformation \(U_K(\omega_K)\)
\cite{poon_law_stochastic_pmd_2008}. The complete two-arm operator is
\(\mathcal{U}_A\otimes\mathcal{U}_B\).

Application to the factorized input state gives
%
\begin{equation}
    |\Psi_{\mathrm{out}}\rangle
    =
    \iint
    d\omega_A\,d\omega_B\,
    F(\omega_A,\omega_B)
    |\omega_A,\omega_B\rangle
    \otimes
    \left[
        U_A(\omega_A)
        \otimes
        U_B(\omega_B)
    \right]
    |\Psi^+\rangle_{\mathrm{pol}}.
    \label{eq:frequency_dependent_spdc_output_state}
\end{equation}

Although the complete evolution remains unitary, the output polarization
state now depends on \((\omega_A,\omega_B)\). Polarization and frequency are
therefore generally no longer separable.

Writing
\(\rho_{AB}^{\mathrm{in}}
=|\Psi^+\rangle\langle\Psi^+|\),
tracing over the unresolved frequencies yields
%
\begin{equation}
    \rho_{AB}^{(\mathrm{pol})}
    =
    \iint
    d\omega_A\,d\omega_B\,
    J(\omega_A,\omega_B)
    \left[
        U_A(\omega_A)
        \otimes
        U_B(\omega_B)
    \right]
    \rho_{AB}^{\mathrm{in}}
    \left[
        U_A(\omega_A)
        \otimes
        U_B(\omega_B)
    \right]^\dagger.
    \label{eq:frequency_averaged_polarization_state}
\end{equation}

This equation is the central propagation model of the chapter. It connects
the joint spectrum of the SPDC source, the frequency-dependent response of
the two fibers, and the reduced polarization state accessible to
polarization measurements.

The resulting map is a continuous random-unitary channel whose weights are
fixed by the physical joint spectrum. If the spectrum factorizes, the two
frequency averages define independent local channels. If the frequencies
are correlated, the transformations remain local for every frequency pair,
but their statistical sampling is correlated between the two arms
\cite{poon_law_stochastic_pmd_2008}.

The next section specifies the physical structure of \(U_K(\omega_K)\)
through the model of a locally uniform birefringent fiber segment.

%========================================================
% Section 6.3 — Single Frequency-Dependent Segment
%========================================================

\section{Single Frequency-Dependent Birefringent Segment}
\label{sec:single_frequency_dependent_segment}

The elementary component of the frequency-dependent fiber model is a locally
uniform birefringent segment. Within the segment, the birefringence axis is
assumed to remain constant, while two orthogonal polarization eigenmodes
accumulate different frequency-dependent phases. This is the local
description underlying first-order PMD models
\cite{foschini_poole_statistical_pmd_1991,
gordon_kogelnik_pmd_2000}.

%--------------------------------------------------------
% Subsection 6.3.1 — Segment Operator and DGD
%--------------------------------------------------------

\subsection{Segment Operator and Differential Group Delay}
\label{subsec:single_segment_operator_and_dgd}

Let \(|s_{j,+}\rangle\) and \(|s_{j,-}\rangle\) denote the orthonormal
polarization eigenmodes of segment \(j\). Their projectors form a complete
basis of the polarization-qubit space. For a lossless segment, the
frequency-dependent polarization operator is
%
\begin{equation}
    U_j(\omega)
    =
    e^{i\phi_{j,+}(\omega)}
    |s_{j,+}\rangle\langle s_{j,+}|
    +
    e^{i\phi_{j,-}(\omega)}
    |s_{j,-}\rangle\langle s_{j,-}|,
    \label{eq:single_segment_jones_operator}
\end{equation}
%
where \(\phi_{j,\pm}(\omega)\) are the phases accumulated by the two
eigenmodes.

The differential phase is
%
\begin{equation}
    \Delta\phi_j(\omega)
    =
    \phi_{j,+}(\omega)
    -
    \phi_{j,-}(\omega).
    \label{eq:single_segment_differential_phase}
\end{equation}
%
The average phase
\(\overline{\phi}_j(\omega)
=[\phi_{j,+}(\omega)+\phi_{j,-}(\omega)]/2\)
is common to both polarization components and therefore does not affect the
reduced polarization state.

Let \(\widehat{\mathbf n}_j\) be the Bloch-sphere axis associated with the
two eigenmodes, so that their projectors are
\((I_2\pm\widehat{\mathbf n}_j\cdot\boldsymbol{\sigma})/2\).
Equation~\eqref{eq:single_segment_jones_operator} can then be written as
%
\begin{equation}
    U_j(\omega)
    =
    e^{i\overline{\phi}_j(\omega)}
    \exp
    \left[
        \frac{i}{2}
        \Delta\phi_j(\omega)
        \widehat{\mathbf n}_j
        \cdot
        \boldsymbol{\sigma}
    \right].
    \label{eq:single_segment_pauli_exponential}
\end{equation}

Apart from its common phase, the segment therefore acts as a polarization
rotation about \(\widehat{\mathbf n}_j\), with a rotation angle that depends
on frequency.

\medskip

Around the carrier frequency \(\omega_0\), the signed differential group
delay is defined by
%
\begin{equation}
    \Delta\tau_j
    =
    \left.
    \frac{d\Delta\phi_j(\omega)}
    {d\omega}
    \right|_{\omega=\omega_0}.
    \label{eq:single_segment_signed_differential_group_delay}
\end{equation}
%
Its sign depends on the labeling of the two eigenmodes, while their physical
temporal separation is
%
\begin{equation}
    \operatorname{DGD}_j
    =
    |\Delta\tau_j|.
    \label{eq:single_segment_dgd}
\end{equation}

In the first-order PMD approximation, the differential phase is linearized
as
%
\begin{equation}
    \Delta\phi_j(\omega)
    \simeq
    \Delta\phi_j(\omega_0)
    +
    (\omega-\omega_0)\Delta\tau_j.
    \label{eq:single_segment_phase_linearization}
\end{equation}
%
This approximation assumes that the birefringence axis and differential
delay remain approximately constant over the photon bandwidth
\cite{gordon_kogelnik_pmd_2000}.

After removing the common phase and the deterministic polarization rotation
at \(\omega_0\), the residual frequency-dependent operator is
%
\begin{equation}
    U_j^{(\mathrm{res})}(\omega)
    \simeq
    \exp
    \left[
        \frac{i}{2}
        (\omega-\omega_0)
        \Delta\tau_j
        \widehat{\mathbf n}_j
        \cdot
        \boldsymbol{\sigma}
    \right].
    \label{eq:single_segment_residual_operator}
\end{equation}

This residual operator contains the part of the segment response responsible
for polarization--frequency coupling. The carrier-frequency rotation is
coherent and can change the polarization basis, but it cannot by itself
reduce purity or entanglement.

%--------------------------------------------------------
% Subsection 6.3.2 — Polarization--Frequency Coupling
%--------------------------------------------------------

\subsection{Polarization--Frequency Coupling}
\label{subsec:single_segment_polarization_frequency_coupling}

Consider a photon with normalized spectral amplitude \(f(\omega)\) and
polarization
\(c_+|s_{j,+}\rangle+c_-|s_{j,-}\rangle\).
In the compensated frame defined above, propagation produces
%
\begin{align}
    |\psi_{\mathrm{out}}\rangle
    \simeq
    \int d\omega\,
    f(\omega)|\omega\rangle
    \otimes
    \Big[
        &
        c_+
        e^{\frac{i}{2}(\omega-\omega_0)\Delta\tau_j}
        |s_{j,+}\rangle
        \nonumber\\
        {}+
        &
        c_-
        e^{-\frac{i}{2}(\omega-\omega_0)\Delta\tau_j}
        |s_{j,-}\rangle
    \Big].
    \label{eq:single_segment_output_wave_packet}
\end{align}

The complete state remains pure, but polarization and frequency are generally
no longer separable. When frequency is not resolved, the off-diagonal
polarization coherence is multiplied by
%
\begin{equation}
    \kappa_j
    =
    \int d\omega\,
    |f(\omega)|^2
    e^{i(\omega-\omega_0)\Delta\tau_j}.
    \label{eq:single_segment_coherence_factor}
\end{equation}

In the segment eigenmode basis, the reduced polarization state becomes
%
\begin{equation}
    \rho_{\mathrm{pol}}^{(\mathrm{out})}
    =
    \begin{pmatrix}
        \rho_{++}^{(\mathrm{in})}
        &
        \kappa_j\rho_{+-}^{(\mathrm{in})}
        \\
        \kappa_j^*\rho_{-+}^{(\mathrm{in})}
        &
        \rho_{--}^{(\mathrm{in})}
    \end{pmatrix}.
    \label{eq:single_segment_reduced_polarization_state}
\end{equation}

The magnitude \(|\kappa_j|\) quantifies the residual overlap of the two
temporally displaced polarization components. For
\(|\kappa_j|=1\), the reduced polarization evolution remains coherent. For
\(|\kappa_j|<1\), tracing over frequency suppresses the polarization
off-diagonal terms and produces effective decoherence
\cite{poon_law_stochastic_pmd_2008}.

The single-segment model thus identifies the elementary quantum mechanism of
PMD: frequency-dependent unitary propagation creates
polarization--frequency correlations, which appear as polarization
decoherence only after the unresolved frequency degree of freedom is traced
out. A nonuniform fiber can then be modeled by concatenating segments with
different birefringence axes and differential delays.
%========================================================
% Section 6.4 — Concatenated Birefringent Fiber
%========================================================

\section{Concatenated Birefringent Fiber Model}
\label{sec:concatenated_birefringent_fiber}

In a real communication fiber, the local birefringence generally varies
along the propagation direction because of geometrical imperfections,
mechanical stress, bending, and environmental perturbations. The fiber is
therefore represented as a sequence of locally uniform segments, each
described by the frequency-dependent operator introduced in the previous
section
\cite{foschini_poole_statistical_pmd_1991,
wai_menyuk_random_birefringence_1996,
gordon_kogelnik_pmd_2000}.

%--------------------------------------------------------
% Subsection 6.4.1 — Ordered Propagation
%--------------------------------------------------------

\subsection{Ordered Propagation Along a Fiber Arm}
\label{subsec:ordered_segment_propagation}

Let \(K\in\{A,B\}\) denote one of the two fiber arms and let \(N_K\) be its
number of segments. If the photon encounters segment \(1\) first and segment
\(N_K\) last, the total operator is
%
\begin{equation}
    U_K(\omega_K)
    =
    U_{K,N_K}(\omega_K)
    \cdots
    U_{K,2}(\omega_K)
    U_{K,1}(\omega_K).
    \label{eq:arm_concatenated_fiber_operator}
\end{equation}

The rightmost operator therefore acts first. This ordering is physically
relevant because segments with different birefringence axes generally do not
commute:
%
\begin{equation}
    U_{K,j}(\omega_K)
    U_{K,k}(\omega_K)
    \neq
    U_{K,k}(\omega_K)
    U_{K,j}(\omega_K).
    \label{eq:segment_operator_noncommutativity}
\end{equation}

Each segment changes the polarization basis in which the following segment
acts. Consequently, the complete transformation depends on both the local
segment parameters and their spatial ordering
\cite{gordon_kogelnik_pmd_2000,
wai_menyuk_random_birefringence_1996}.

A useful limiting case occurs when all segments share the same oriented
birefringence axis \(\widehat{\mathbf n}_K\). The operators then commute and
the product reduces to
%
\begin{align}
    U_K(\omega_K)
    =
    &
    \exp
    \left[
        i
        \sum_{j=1}^{N_K}
        \overline{\phi}_{K,j}(\omega_K)
    \right]
    \nonumber\\
    &\times
    \exp
    \left[
        \frac{i}{2}
        \left(
            \sum_{j=1}^{N_K}
            \Delta\phi_{K,j}(\omega_K)
        \right)
        \widehat{\mathbf n}_K
        \cdot
        \boldsymbol{\sigma}
    \right].
    \label{eq:aligned_segments_total_operator}
\end{align}

In this special case, the differential phases and signed differential delays
add directly, and the fiber is equivalent to a single birefringent segment.
For non-aligned axes, no corresponding scalar addition rule exists: the
ordered product generates effective polarization modes and a total DGD that
must be extracted from the complete operator.

%--------------------------------------------------------
% Subsection 6.4.2 — Fiber Realizations
%--------------------------------------------------------

\subsection{Fixed and Statistical Fiber Realizations}
\label{subsec:fixed_and_statistical_fiber_realizations}

A fiber realization is specified by the birefringence axes and differential
phases of its segments. These quantities may be prescribed deterministically
or sampled from a statistical model describing spatially varying
birefringence.

For any fixed realization and resolved frequency, every segment operator is
unitary and so is their ordered product. Spatial randomness inside the model
does not therefore produce mixedness by itself. The reduced polarization
state becomes mixed only when frequency is traced out or when measurements
average over different fiber configurations
\cite{wai_menyuk_random_birefringence_1996,
poon_law_stochastic_pmd_2008}.

The two fiber arms are constructed independently through
Eq.~\eqref{eq:arm_concatenated_fiber_operator}. Their resulting operators
\(U_A(\omega_A)\) and \(U_B(\omega_B)\) are then used in the two-photon
spectral propagation model developed in the previous section.

The concatenated construction thus supplies the complete
frequency-dependent response of each fiber arm. The next section extracts
from this response the effective principal states of polarization and the
total differential group delay.

%========================================================
% Section 6.5 — Effective PSPs and Total DGD
%========================================================

\section{Effective Principal States of Polarization and Total DGD}
\label{sec:effective_psp_and_total_dgd}

The ordered product introduced in the previous section provides the complete
frequency-dependent operator \(U_K(\omega_K)\) of fiber arm \(K\). The local
eigenmodes and delays of its individual segments do not, however, directly
represent the first-order response of the complete fiber. Effective
principal states and a total differential group delay must be extracted from
the frequency dependence of the concatenated operator
\cite{gordon_kogelnik_pmd_2000}.

%--------------------------------------------------------
% Subsection 6.5.1 — PMD Generator and PSPs
%--------------------------------------------------------

\subsection{PMD Generator and Principal States}
\label{subsec:complete_fiber_pmd_generator}

The complete fiber operator is separated into a polarization-independent
phase and a polarization transformation:
%
\begin{equation}
    U_K(\omega_K)
    =
    e^{i\Phi_K(\omega_K)}
    W_K(\omega_K),
    \qquad
    W_K(\omega_K)\in\mathrm{SU}(2).
    \label{eq:total_operator_recalled_for_pmd}
\end{equation}

The derivative of the common phase describes an average delay shared by both
polarization components and does not affect the reduced polarization state.
The polarization-dependent first-order response is instead described by the
input PMD generator
%
\begin{equation}
    G_K^{(\mathrm{in})}(\omega_K)
    =
    -i
    W_K^\dagger(\omega_K)
    \frac{
        \partial W_K(\omega_K)
    }{
        \partial\omega_K
    }.
    \label{eq:input_pmd_generator}
\end{equation}

Because \(W_K(\omega_K)\) is unitary with unit determinant,
\(G_K^{(\mathrm{in})}(\omega_K)\) is Hermitian and traceless. Its
eigenvectors at the carrier frequency \(\omega_{0,K}\) define the input
principal states of polarization:
%
\begin{equation}
    G_K^{(\mathrm{in})}(\omega_{0,K})
    |p_{K,\pm}^{(\mathrm{in})}\rangle
    =
    \pm
    \frac{\tau_{\mathrm{DGD},K}}{2}
    |p_{K,\pm}^{(\mathrm{in})}\rangle.
    \label{eq:complete_fiber_psp_eigenproblem}
\end{equation}

The eigenvalue separation is the non-negative total differential group delay
\(\tau_{\mathrm{DGD},K}\). Since the generator is Hermitian, the two PSPs
form an orthonormal polarization basis.

The corresponding output PSPs are
%
\begin{equation}
    |p_{K,\pm}^{(\mathrm{out})}\rangle
    =
    W_K(\omega_{0,K})
    |p_{K,\pm}^{(\mathrm{in})}\rangle.
    \label{eq:output_psps_from_input_psps}
\end{equation}

To first order around the carrier frequency, propagation of an input PSP
satisfies
%
\begin{equation}
    W_K(\omega_K)
    |p_{K,\pm}^{(\mathrm{in})}\rangle
    \simeq
    e^{
        \pm
        i(\omega_K-\omega_{0,K})
        \tau_{\mathrm{DGD},K}/2
    }
    |p_{K,\pm}^{(\mathrm{out})}\rangle.
    \label{eq:first_order_psp_propagation}
\end{equation}

Thus, all spectral components launched in one PSP emerge with the same
output polarization to first order. They differ only by a
frequency-dependent phase corresponding to an early or late temporal
displacement. A superposition of the two PSPs therefore becomes correlated
with arrival time, providing the physical mechanism through which PMD can
reduce polarization coherence and entanglement after the temporal or
spectral degree of freedom is traced out
\cite{poon_law_stochastic_pmd_2008,
riccardi_polarization_entanglement_distribution_2021}.

The PSPs of the complete fiber should not be confused with the local
eigenmodes of an individual segment. Local eigenmodes diagonalize one
segment operator, whereas the PSPs diagonalize the first-order frequency
response of the complete concatenated fiber.

%--------------------------------------------------------
% Subsection 6.5.2 — PMD Vector and Numerical Extraction
%--------------------------------------------------------

\subsection{PMD Vector and Numerical Extraction}
\label{subsec:numerical_psp_dgd_extraction}

Since the PMD generator is a traceless Hermitian \(2\times2\) matrix, it can
be represented in the Pauli basis as
%
\begin{equation}
    G_K^{(\mathrm{in})}(\omega_{0,K})
    =
    \frac{1}{2}
    \boldsymbol{\tau}_K
    \cdot
    \boldsymbol{\sigma},
    \qquad
    \tau_{\mathrm{DGD},K}
    =
    \lVert\boldsymbol{\tau}_K\rVert.
    \label{eq:complete_pmd_vector_representation}
\end{equation}

The direction of \(\boldsymbol{\tau}_K\) defines the PSP axis on the Bloch
sphere, while its magnitude gives the total DGD. The vector already includes
the ordered accumulation of all local segment contributions; these
contributions need not be reconstructed separately once the complete
operator \(W_K(\omega_K)\) is available.

Numerically, its frequency derivative can be evaluated around
\(\omega_{0,K}\) using a centered finite difference:
%
\begin{equation}
    \left.
    \frac{
        \partial W_K
    }{
        \partial\omega_K
    }
    \right|_{\omega_{0,K}}
    \simeq
    \frac{
        W_K(\omega_{0,K}+\Delta\omega_K)
        -
        W_K(\omega_{0,K}-\Delta\omega_K)
    }{
        2\Delta\omega_K
    }.
    \label{eq:finite_difference_total_operator_derivative}
\end{equation}

Substitution into Eq.~\eqref{eq:input_pmd_generator}, followed by
diagonalization, gives the numerical PSPs and total DGD. The result should
be verified by reducing \(\Delta\omega_K\) until the extracted quantities
are stable. A continuous phase convention for \(W_K(\omega_K)\) must be
maintained across the frequency samples so that the derivative represents
the physical variation of the operator rather than an arbitrary phase
discontinuity.

These quantities provide a compact first-order description of each complete
fiber arm. Their effect on the entangled photon pair is nevertheless
determined jointly by the DGD, the PSP orientations, and the two-photon
spectrum through the reduced-state propagation model.

%========================================================
% Section 6.6 — Numerical Construction of the
%               Segmented-Fiber Model
%========================================================

\section{Numerical Construction of the Segmented-Fiber Model}
\label{sec:numerical_segmented_fiber_construction}

The numerical model evaluates the frequency-dependent fiber operators
introduced in the preceding sections and approximates the spectral trace by
a finite quadrature. For each fiber realization, the calculation therefore
requires a discrete joint spectrum and the ordered segment product evaluated
at every frequency node.

The spatial parameters defining a fiber realization, such as the local
birefringence axes, differential phases, and signed differential delays, are
kept fixed during the spectral propagation. If stochastic fibers are
considered, their probability distributions must be specified as part of
the numerical experiment
\cite{foschini_poole_statistical_pmd_1991,
gordon_kogelnik_pmd_2000}.

%--------------------------------------------------------
% Subsection 6.6.1 — Spectral Discretization
%--------------------------------------------------------

\subsection{Discretization of the Joint Spectrum}
\label{subsec:joint_spectrum_discretization}

Let
\(\{\omega_{A,m}\}_{m=1}^{M_A}\) and
\(\{\omega_{B,n}\}_{n=1}^{M_B}\) denote the frequency nodes used for the two
photons, with corresponding quadrature weights \(a_m\) and \(b_n\).
The joint spectral intensity \(J(\omega_A,\omega_B)\), defined in
Eq.~\eqref{eq:joint_spectral_intensity}, is represented by the normalized
discrete weights
%
\begin{equation}
    Q_{mn}
    =
    \frac{
        a_m b_n
        J(\omega_{A,m},\omega_{B,n})
    }{
        \displaystyle
        \sum_{m'=1}^{M_A}
        \sum_{n'=1}^{M_B}
        a_{m'} b_{n'}
        J(\omega_{A,m'},\omega_{B,n'})
    },
    \qquad
    Q_{mn}\geq 0,
    \qquad
    \sum_{m,n}Q_{mn}=1.
    \label{eq:normalized_discrete_joint_spectral_weights}
\end{equation}

The explicit normalization compensates for the finite spectral interval and
quadrature resolution. Numerical convergence must consequently be checked
with respect to both the number of frequency nodes and the spectral interval
included in the calculation.

%--------------------------------------------------------
% Subsection 6.6.2 — Discrete Reduced State
%--------------------------------------------------------

\subsection{Discrete Reduced Polarization State}
\label{subsec:ordered_numerical_spectral_propagation}

For a fixed fiber realization \(r\), the segment operators are evaluated at
each frequency node and multiplied in the propagation order specified by
Eq.~\eqref{eq:arm_concatenated_fiber_operator}. After removing the
polarization-independent phase according to
Eq.~\eqref{eq:total_operator_recalled_for_pmd}, the resulting polarization
operators are denoted by
\(W_A^{(r)}(\omega_{A,m})\) and
\(W_B^{(r)}(\omega_{B,n})\).

For the frequency pair
\((\omega_{A,m},\omega_{B,n})\), the corresponding two-photon operator is
%
\begin{equation}
    \mathcal{W}_{mn}^{(r)}
    =
    W_A^{(r)}(\omega_{A,m})
    \otimes
    W_B^{(r)}(\omega_{B,n}).
    \label{eq:discrete_two_photon_polarization_operator}
\end{equation}

The discrete approximation of the reduced polarization state in
Eq.~\eqref{eq:frequency_averaged_polarization_state} is then
%
\begin{equation}
    \rho_{AB}^{(\mathrm{pol},r)}
    \simeq
    \sum_{m=1}^{M_A}
    \sum_{n=1}^{M_B}
    Q_{mn}
    \mathcal{W}_{mn}^{(r)}
    \rho_{AB}^{(\mathrm{in})}
    \mathcal{W}_{mn}^{(r)\dagger}.
    \label{eq:discrete_frequency_averaged_polarization_state}
\end{equation}

Equation~\eqref{eq:discrete_frequency_averaged_polarization_state}
is a weighted sum of the polarization density operators associated with the
different frequency pairs. The Jones operators themselves must not be
averaged before acting on the input state, since that procedure would
describe a different physical transformation.

The corresponding Kraus operators are
%
\begin{equation}
    K_{mn}^{(r)}
    =
    \sqrt{Q_{mn}}\,
    \mathcal{W}_{mn}^{(r)},
    \qquad
    \sum_{m,n}
    K_{mn}^{(r)\dagger}K_{mn}^{(r)}
    =
    I_4.
    \label{eq:discrete_spectral_kraus_operators}
\end{equation}

The discretized transformation is therefore completely positive and trace
preserving. Its effective non-unitarity results from tracing over the
unresolved frequency degree of freedom, while propagation at each frequency
remains locally unitary
\cite{nielsen_chuang_qcqi,
poon_law_stochastic_pmd_2008,
riccardi_polarization_entanglement_distribution_2021}.

%--------------------------------------------------------
% Subsection 6.6.3 — Spectral Correlations
%--------------------------------------------------------

\subsection{Factorized and Correlated Spectra}
\label{subsec:factorized_correlated_spectral_sampling}

If the discrete spectral distribution factorizes as
%
\begin{equation}
    Q_{mn}
    =
    q_{A,m}q_{B,n},
    \label{eq:factorized_discrete_joint_weights}
\end{equation}
%
the spectral trace is equivalent to the application of two independently
averaged local polarization channels.

For a correlated joint spectrum, the transformation at every frequency pair
remains local because
\(\mathcal{W}_{mn}^{(r)}\) is a tensor product. Nevertheless, the local
operators sampled in the two arms are statistically correlated through
\(Q_{mn}\). The resulting reduced channel therefore cannot generally be
written as the tensor product of two independent averaged channels.

This distinction separates local PMD-induced decoherence from correlations
in the effective polarization noise inherited from the joint spectrum of
the photon pair.

%--------------------------------------------------------
% Subsection 6.6.4 — Fiber Realizations
%--------------------------------------------------------

\subsection{Averaging over Fiber Realizations}
\label{subsec:numerical_fiber_realization_averaging}

Equation~\eqref{eq:discrete_frequency_averaged_polarization_state}
describes the spectral trace for one fixed pair of fibers. If the
birefringence configuration changes during the measurement, or if an
ensemble of fibers is considered, a second statistical average is required:
%
\begin{equation}
    \overline{\rho}_{AB}^{(\mathrm{pol})}
    =
    \sum_{r=1}^{N_{\mathrm{R}}}
    \pi_r
    \rho_{AB}^{(\mathrm{pol},r)},
    \qquad
    \pi_r\geq0,
    \qquad
    \sum_{r=1}^{N_{\mathrm{R}}}\pi_r=1.
    \label{eq:numerical_fiber_realization_average}
\end{equation}

For equally weighted Monte Carlo realizations,
\(\pi_r=1/N_{\mathrm{R}}\). The spectral trace must first be completed for
each realization; only afterwards are the resulting density operators
averaged. This preserves the distinction between unresolved frequencies
within one photon wave packet and classical uncertainty about the fiber
configuration.

When the realization is not experimentally resolved, purity, fidelity,
concurrence, and Bell parameters must be evaluated from
\(\overline{\rho}_{AB}^{(\mathrm{pol})}\). Averaging their
realization-dependent values instead characterizes the statistical
distribution of individual fibers and represents a different physical
question.

%========================================================
% Section 6.7 — Limiting Cases and Validation
%========================================================

\section{Limiting Cases and Numerical Validation}
\label{sec:segmented_fiber_validation}

The numerical model must reproduce analytically controlled limits before it
is applied to general concatenated fibers. The most relevant tests concern
the recovery of the single-segment model, the aligned-axis limit, the
frequency-independent limit, and convergence of the reduced polarization
state.

%--------------------------------------------------------
% Analytically Controlled Limits
%--------------------------------------------------------

\subsection{Analytically Controlled Limits}

For a fiber containing a single segment, the total operator must reduce to
the segment operator introduced in Section~6.3. The extracted principal
states must coincide, up to their ordering and arbitrary global phases, with
the local birefringence eigenmodes. The corresponding DGD must equal the
magnitude of the signed differential delay \(\Delta\tau_{K,1}\).

After spectral tracing, the polarization coherence must reproduce the factor
defined in Eq.~\eqref{eq:single_segment_coherence_factor}. For a Gaussian
spectral probability distribution with standard deviation
\(\sigma_{\omega,K}\), its magnitude is
%
\begin{equation}
    \left|\kappa_{K,1}\right|
    =
    \exp
    \left[
        -
        \frac{1}{2}
        \sigma_{\omega,K}^{2}
        \Delta\tau_{K,1}^{2}
    \right].
    \label{eq:single_segment_gaussian_validation_target}
\end{equation}

The numerical spectral average must converge toward
Eq.~\eqref{eq:single_segment_gaussian_validation_target} as the frequency
grid is refined and the integration interval is enlarged
\cite{poon_law_stochastic_pmd_2008}.

A second analytical limit is obtained when all segment axes are aligned.
The ordered product must then reproduce the commuting operator given in
Eq.~\eqref{eq:aligned_segments_total_operator}, and the signed differential
delays must add directly. This test verifies both the multiplication order
and the adopted sign convention.

Finally, in the monochromatic limit, or whenever all segment operators are
frequency independent, the discrete spectral average reduces to a single
local-unitary transformation:
%
\[
    \rho_{AB}^{(\mathrm{pol})}
    =
    \mathcal{W}
    \rho_{AB}^{(\mathrm{in})}
    \mathcal{W}^{\dagger}.
\]
%
A pure input state must therefore remain pure, and its concurrence must be
preserved. Fidelity with a fixed Bell state may change because the local
unitaries can coherently rotate the polarization basis. This limit
distinguishes coherent birefringent rotation from PMD-induced polarization
decoherence.

%--------------------------------------------------------
% Numerical Convergence
%--------------------------------------------------------

\subsection{Numerical Convergence}

The frequency nodes must cover the significant support of the joint
spectrum. Convergence must therefore be checked with respect to both grid
resolution and integration interval.

Let \(\rho_{AB}^{(M_A,M_B)}\) denote the reduced state obtained with
\(M_A\) and \(M_B\) frequency nodes. A convenient grid-convergence indicator
is
%
\begin{equation}
    \epsilon_{\rho}^{(\mathrm{grid})}
    =
    \left\|
        \rho_{AB}^{(2M_A,2M_B)}
        -
        \rho_{AB}^{(M_A,M_B)}
    \right\|_{\mathrm{F}}.
    \label{eq:density_matrix_grid_convergence_error}
\end{equation}

The predicted purity, fidelity, concurrence, and Bell parameter should also
remain stable under the same refinement. If the PMD generator is evaluated
through finite differences, its DGD must additionally be stable when the
frequency step is reduced.

For every calculation, the following numerical properties must be verified:

\begin{itemize}
    \item unitarity of each fixed-frequency fiber operator;
    \item normalization of the discrete spectral weights;
    \item Hermiticity and unit trace of the reduced density operator;
    \item positivity of the reduced density operator within numerical
    precision;
    \item recovery of the single-segment, aligned-axis, and
    frequency-independent limits.
\end{itemize}

These tests are sufficient to validate the construction of the quantum state
without introducing a separate statistical theory of fiber ensembles.

%========================================================
% Section 6.8 — Chapter Summary
%========================================================

\section{Chapter Summary}
\label{sec:chapter_frequency_dependent_fiber_summary}

This chapter extended the effective polarization models of the preceding
chapter by restoring the spectral degree of freedom of the photons. Each
fiber arm was therefore described by a frequency-dependent Jones operator
rather than by a single polarization gate.

Propagation remains unitary in the complete
frequency--polarization Hilbert space. Frequency-dependent birefringence can,
however, correlate polarization with frequency and time of arrival. When
these degrees of freedom are not resolved, their partial trace produces an
effective mixed polarization state.

A locally uniform birefringent segment was introduced as the elementary
component of the model. Its orthogonal polarization eigenmodes acquire
different spectral phases and group delays. A spatially nonuniform fiber was
then represented as an ordered product of such segments. Because rotations
around different birefringence axes do not generally commute, the resulting
transformation depends on their propagation order.

The frequency derivative of the complete fiber operator defines its
effective principal states of polarization and total differential group
delay. These quantities characterize the first-order frequency sensitivity
of the complete fiber without requiring the local segment delays to be added
as scalars.

Finally, the continuous spectral trace was converted into a normalized
quadrature over the joint spectrum. The resulting reduced polarization state
defines a random-unitary quantum channel whose local transformations may be
statistically correlated through the photon-pair spectrum. This framework
provides the basis for evaluating how PMD modifies the purity, fidelity,
entanglement, and nonlocal correlations of polarization-entangled photon
pairs.